\begin{textsample}{NEG Dim 8 – global\_south\_2023\_10 – Score -1.00 – global\_south\_2023\_10/103031488.txt}  \label{ex:f8_neg_090}
UN reso on Gaza a ' dishonor ' to memory of 4 OFWs -- DFA Ambassador Antonio Lagdameo, Philippine Permanent Mission Representative to the UN ( right ), explains why the Philippines abstained on the Jordan-sponsored resolution before the UNGA emergency special session in New York Friday.
To his left is Deputy Permanent Mission Representative Ariel Pe? aranda.
THE Philippine government fully supports the call for a ceasefire in Gaza Strip, but was forced to abstain from an Arab-sponsored resolution in the United Nations General Assembly because it fell short of condemning the Hamas attacks that killed 1,400 people including four \textbf{Filipinos}.
The UNGA passed a resolution calling for an"immediate, durable and sustained humanitarian truce"between Israeli forces and Hamas militants in Gaza Strip."The Resolution as worded contained many elements that we do support, such as the urgency for access to humanitarian assistance to Gaza and for all parties to adhere to international humanitarian law."But for the Philippines to have voted for a Resolution which simply condemns Israeli actions without mentioning how this situation was precipitated @ @ @ @ @ @ @ @ @ @ of \textbf{Filipinos}, would be to dishonor the memory of our murdered citizens.
That is why we voted to abstain,"DFA Undersecretary for Migrant Workers Eduardo de Vega said.
UNGA resolution The Jordan-initiated resolution calls for"immediate, durable and sustained humanitarian truce leading to a cessation of hostilities."A total of 121 countries or a two-thirds majority of the UNGA voted in favor of it, while 14 including the United States and Israel voted against it.
The Philippines, along with 44 other countries, abstained.
The UNGA resolution also"demands"for all parties to protect civilians, humanitarian personnel and facilitate humanitarian access for essential supplies and services.
Civilians who are"illegally being held captive"should be immediately and unconditionally be released.
DFA statement In a separate statement, the DFA said Canada introduced an amendment to the Jordan resolution.
Canada proposed to include mentioning the October 7"terrorist attacks,"which killed four \textbf{Filipino} caregivers.
Two other OFWs @ @ @ @ @ @ @ @ @ @."In regard to this Philippine interest, we supported Canada ' s proposal to achieve more balance in the draft, with a factual reference to, and condemnation of, the 7 October terrorist attacks by Hamas that killed many innocent civilians including \textbf{Filipinos} working and living in Israel."Canada ' s proposal was supported by 88 states, but we regret that it fell short of 8 more votes that would have seen this critical element, which is important to the Philippines as to other countries, reflected in an important UN resolution,"DFA spokesperson Ma.
Teresita Daza said.
Daza stressed that the Philippines has been in"solidarity with the global community"in calling for swift action to stop the"scale of human suffering"affecting both Israel and Gaza Strip.
She noted that the Philippine position has been consistent with the Asean and Asean-Gulf Cooperation Council statements, and during the speech made by the Philippine Mission to the UN in New York during the UN Security Council open debate. @ @ @ @ @ @ @ @ @ @ shares with the global community -- as conveyed in our statements -- including grave concerns on acts of violence against all civilians, calls for all parties to adhere to international humanitarian law, the urgency of access for humanitarian assistance to address the large-scale needs in Gaza, and the immediate and unconditional release of innocent civilians being held captive.
It also affirms long-standing support for a two-State solution with a safe and independent Palestine and a secure Israel living in peace,"she said.
The DFA commits that the Philippines will continue to support the UN, the UN Security Council, the UN humanitarian system and the global community"to decisively put a stop to the alarming deaths and suffering in Gaza and Israel, assist those in need of humanitarian help, and restore peace and normalcy to the lives of millions of people affected by this crisis."

% matched lemmas: filipino
\end{textsample}
